\documentclass[UTF8,a4paper,11pt,fontset=none]{ctexart}
\usepackage[left=24mm,right=24mm,top=24mm,bottom=24mm,headheight=16pt]{geometry}
\usepackage{fontspec}
% Open fonts distributed with TeX Live keep the source portable.
\setmainfont{texgyretermes-regular.otf}[BoldFont=texgyretermes-bold.otf,ItalicFont=texgyretermes-italic.otf,BoldItalicFont=texgyretermes-bolditalic.otf]
\setsansfont{texgyreheros-regular.otf}[BoldFont=texgyreheros-bold.otf]
\setmonofont{texgyrecursor-regular.otf}[Scale=MatchLowercase]
\setCJKmainfont{FandolSong-Regular.otf}[BoldFont=FandolSong-Bold.otf,ItalicFont=FandolKai-Regular.otf]
\setCJKsansfont{FandolHei-Regular.otf}[BoldFont=FandolHei-Bold.otf]
\setCJKmonofont{FandolFang-Regular.otf}
\usepackage{amsmath,amssymb,booktabs,longtable,array,seqsplit}
\usepackage{xcolor,graphicx,tikz,caption}
\usetikzlibrary{arrows.meta,positioning}
\usepackage{fvextra,needspace}
\fvset{fontsize=\footnotesize,baselinestretch=1,breaklines,breakanywhere,breaksymbolleft={},breaksymbolright={},breakanywheresymbolpre={}}
\definecolor{reportblue}{HTML}{173B57}
\definecolor{reportorange}{HTML}{BF671A}
\usepackage{hyperref}
\hypersetup{colorlinks=true,urlcolor=reportblue,linkcolor=black,unicode=true,hypertexnames=false,pdftitle={香山昆明湖V2自定义向量点积指令技术报告},pdfsubject={设计、验证与性能评估},pdfauthor={啊对对对对}}
\urlstyle{same}
\usepackage{fancyhdr}
\pagestyle{fancy}\fancyhf{}
\fancyhead[L]{\small 香山昆明湖 V2 · 自定义向量点积指令}
\fancyhead[R]{\small 技术报告}
\fancyfoot[C]{\thepage}
\ctexset{section={format=\Large\bfseries\rmfamily\color{black}},subsection={format=\large\bfseries\rmfamily\color{black}},subsubsection={format=\normalsize\bfseries\rmfamily\color{black}}}
\setlength{\parindent}{2em}
\setlength{\parskip}{3pt}
\linespread{1.22}
\setlength{\LTpre}{8pt}\setlength{\LTpost}{8pt}
\renewcommand{\arraystretch}{1.22}
\emergencystretch=3em
\setcounter{tocdepth}{2}
\setcounter{secnumdepth}{2}
\newcommand{\reportteam}{啊对对对对}
\newcommand{\reportcover}[1]{%
\begin{titlepage}\thispagestyle{empty}\centering
\vspace*{17mm}
{\Large\color{black} 香山昆明湖 V2\par}\vspace{9mm}
{\Huge\bfseries\color{black} 自定义向量点积指令\par}\vspace{9mm}
{\LARGE\bfseries #1\par}\vspace{7mm}
{\large 设计、验证与性能评估\par}\vspace{12mm}
\vfill
{\Large 参赛队伍：\reportteam\par}\vspace{10mm}
{\large 三阶段工程技术报告\par}\vspace{5mm}
{\normalsize 2026年11月8日修订\par}\vspace{13mm}
\end{titlepage}}
\newcommand{\reportcontents}{\pagenumbering{Roman}\tableofcontents\clearpage\pagenumbering{arabic}}
\newcommand{\reporttitle}[1]{%
\clearpage\setcounter{section}{0}\setcounter{table}{0}
\phantomsection\addcontentsline{toc}{part}{#1}
\begin{center}{\LARGE\bfseries\color{black} #1\par}\end{center}\vspace{6pt}}
\newcommand{\stagefigureprefix}{1}
\renewcommand{\thesection}{\stagefigureprefix.\arabic{section}}
\renewcommand{\thefigure}{\stagefigureprefix-\arabic{figure}}
\renewcommand{\thetable}{\stagefigureprefix-\arabic{table}}
\captionsetup[figure]{font=small,labelsep=quad,hypcap=false,justification=centering,singlelinecheck=false,position=bottom}
\captionsetup[table]{font=small,labelsep=quad,justification=centering,singlelinecheck=false,position=top}
\setlength{\LTcapwidth}{\textwidth}
% 2026-11-07 统一风格插图：160 mm 画布、1:1 字号（宋体 + Times），源文件见上级目录“源文件”。
% 未列出的 ID（截图与波形）沿用原有素材。
\newcommand{\DeclareReportFigureOverride}[2]{%
\expandafter\newcommand\csname reportfigureasset@#1\endcsname{#2}}
\DeclareReportFigureOverride{S1-01}{figures/unified-20261107/S1-01.pdf}
\DeclareReportFigureOverride{S1-02}{figures/unified-20261107/S1-02.pdf}
\DeclareReportFigureOverride{S1-06}{figures/unified-20261107/S1-06.pdf}
\DeclareReportFigureOverride{S1-07}{figures/unified-20261107/S1-07.pdf}
\DeclareReportFigureOverride{S2-01}{figures/unified-20261107/S2-01.pdf}
\DeclareReportFigureOverride{S2-02}{figures/unified-20261107/S2-02.pdf}
\DeclareReportFigureOverride{S2-03}{figures/unified-20261107/S2-03.pdf}
\DeclareReportFigureOverride{S2-04}{figures/unified-20261107/S2-04.pdf}
\DeclareReportFigureOverride{S2-05}{figures/unified-20261107/S2-05.pdf}
\DeclareReportFigureOverride{S2-06}{figures/unified-20261107/S2-06.pdf}
\DeclareReportFigureOverride{S2-07}{figures/unified-20261107/S2-07.pdf}
\DeclareReportFigureOverride{S2-08}{figures/unified-20261107/S2-08.pdf}
\DeclareReportFigureOverride{S2-09}{figures/unified-20261107/S2-09.pdf}
\DeclareReportFigureOverride{S2-10}{figures/unified-20261107/S2-10.pdf}
\DeclareReportFigureOverride{S3-01}{figures/unified-20261107/S3-01.pdf}
\DeclareReportFigureOverride{S3-02}{figures/unified-20261107/S3-02.pdf}
\DeclareReportFigureOverride{S3-04}{figures/unified-20261107/S3-04.pdf}
\DeclareReportFigureOverride{S3-05}{figures/unified-20261107/S3-05.pdf}
\DeclareReportFigureOverride{S3-06}{figures/unified-20261107/S3-06.pdf}
\DeclareReportFigureOverride{S3-07}{figures/unified-20261107/S3-07.pdf}
\DeclareReportFigureOverride{S3-08}{figures/unified-20261107/S3-08.pdf}
\DeclareReportFigureOverride{S3-09}{figures/unified-20261107/S3-09.pdf}
% 2026-11-08 截图组图：英文子标题改为中文，截图内容不变（源文件/screenshot_figs.py）。
\DeclareReportFigureOverride{S1-03}{figures/unified-20261107/S1-03.pdf}
\DeclareReportFigureOverride{S1-04}{figures/unified-20261107/S1-04.pdf}
\DeclareReportFigureOverride{S1-05}{figures/unified-20261107/S1-05.pdf}
\DeclareReportFigureOverride{SUP-S2-ARITH}{figures/unified-20261107/SUP-S2-ARITH.pdf}
\DeclareReportFigureOverride{SUP-S2-STATE}{figures/unified-20261107/SUP-S2-STATE.pdf}
\DeclareReportFigureOverride{SUP-S2-MODEL}{figures/unified-20261107/SUP-S2-MODEL.pdf}
\DeclareReportFigureOverride{SUP-S2-OVERLAP}{figures/unified-20261107/SUP-S2-OVERLAP.pdf}
\DeclareReportFigureOverride{S2-11}{figures/unified-20261107/S2-11.pdf}

% 图表为浮动体：当前页放不下时移到下一页顶部，由后续正文填补本页；阶段之间由 \clearpage 分隔。
\usepackage{placeins}
\renewcommand{\topfraction}{0.9}\renewcommand{\bottomfraction}{0.6}
\renewcommand{\textfraction}{0.08}\renewcommand{\floatpagefraction}{0.75}
\setcounter{topnumber}{3}\setcounter{bottomnumber}{2}\setcounter{totalnumber}{4}
\setlength{\textfloatsep}{12pt plus 4pt minus 4pt}\setlength{\floatsep}{10pt plus 3pt minus 3pt}
\setlength{\intextsep}{10pt plus 3pt minus 3pt}
\newcommand{\reportfigure}[5]{%
\begin{figure}[!htbp]\centering
\ifcsname reportfigureasset@#1\endcsname
\includegraphics[width=\linewidth,height=#3,keepaspectratio]{\csname reportfigureasset@#1\endcsname}%
\else
\IfFileExists{figures/updated-20261106/#1.pdf}{\includegraphics[width=\linewidth,height=#3,keepaspectratio]{figures/updated-20261106/#1.pdf}}{%
\IfFileExists{figures/#1.pdf}{\includegraphics[width=\linewidth,height=#3,keepaspectratio]{figures/#1.pdf}}{%
\IfFileExists{figures/#1.png}{\includegraphics[width=\linewidth,height=#3,keepaspectratio]{figures/#1.png}}{%
\IfFileExists{figures/#1.jpg}{\includegraphics[width=\linewidth,height=#3,keepaspectratio]{figures/#1.jpg}}{%
\PackageError{report}{Missing figure #1}{Restore the referenced image before compiling.}}}}}%
\fi
\caption{#2}\label{fig:#1}%
% 第4、5参数为作者备注与来源说明，只保留在源码中，不在图下排版；图的解释写在正文。
\end{figure}}
\newcommand{\evidencefigure}[4]{%
\begin{figure}[!htbp]\centering
\ifcsname reportfigureasset@#1\endcsname
\includegraphics[width=\linewidth,height=#3,keepaspectratio]{\csname reportfigureasset@#1\endcsname}%
\else
\includegraphics[width=\linewidth,height=#3,keepaspectratio]{figures/#1.pdf}%
\fi
\caption{#2}\label{fig:#1}%
% 第4参数为作者备注，不在图下排版。
\end{figure}}
